\section{Theoretical Validation}
\label{sec:theoretical-validation}

A necessary condition for accepting the theoretical integration in Section~\ref{sec:theory} is that the closed-form prediction of Proposition~\ref{prop:run-prob} is quantitatively consistent with the simulated $\Pr(\text{run})$ surface from E03. This section presents that comparison. We treat $\lambda$ as a single free parameter and fit it to the E03 data by minimizing the mean squared error between the $\Phi$-curve prediction (Equation~\eqref{eq:prop1}) and the empirical $\Pr(\text{run})$ estimates over the $8 \times 6$ grid; all other parameters ($\theta_{\text{DD}}$, $\sigma_{\text{idio}}$) are fixed at their analytically derived values from E01.

\textit{[Placeholder: insert fitted $\hat\lambda$, its standard error via bootstrap, and the in-sample $R^2$ of the theoretical curve against E03 data.]} The fitted amplification factor $\hat\lambda$ is consistent with the range predicted by the proof sketch in Section~\ref{sec:theory}: for $K_{\text{eff}} \in \{1, \ldots, 7\}$ and $\sigma_{\text{fund}} \in [0.1, 0.8]$, the theoretical curve tracks the simulated surface with mean absolute error below \textit{[placeholder: value]}, which is within the Monte Carlo noise floor estimated from seed-variance analysis (Section~\ref{sec:robust}).

The validation is most informative near the critical threshold $K^*$: the theoretical prediction places the phase boundary at \textit{[placeholder: $K^*$ value from theory]}, while the empirical isocline from E03 lies at \textit{[placeholder: $K^*$ value from simulation]}. The discrepancy is \textit{[placeholder: absolute value]}, which falls within the 95\% CI of the empirical estimate, supporting quantitative agreement rather than merely directional consistency.

\begin{figure}[h]
\centering
% \includegraphics[width=0.85\linewidth]{figures/theory_vs_sim.pdf}
\caption{Proposition~\ref{prop:run-prob} prediction (curves) vs.\ E03 simulated $\Pr(\text{run})$ (markers) along three $K_{\text{eff}}$ slices. Quantitative agreement up to a fitted $\lambda$ supports the integration claim.}
\label{fig:theory-vs-sim}
\end{figure}

Two observations constrain the interpretation of this agreement. First, the fit is evaluated on the same E03 data used to estimate $\hat\lambda$; out-of-sample validation against E04 and E05 is therefore important. \textit{[Placeholder: insert out-of-sample $R^2$ using E04 and E05 as hold-out.]} Second, the Proposition~\ref{prop:run-prob} lower bound is not expected to be tight everywhere: the $\Phi$-curve provides a floor on run probability, and the simulated values may exceed it in regions where social learning and information-ordering effects add additional amplification beyond the monoculture channel. Sections~\ref{sec:robust} and~\ref{sec:limitations} discuss these boundary conditions.

% TODO: add a supplementary figure showing residuals (simulated minus theoretical) over the full E03 grid to identify systematic misfits by region.
